\subsection{Gjeneratorët pseudorastësorë si sisteme dinamike diskrete}
Një gjenerator pseudorastësor është një algoritëm determinist me gjendje të fundme. Gjeneratori kongruencial linear,
\begin{equation}
 x_{n+1}=(ax_n+c)\bmod m,
 \qquad u_n=\frac{x_n}{m},
\end{equation}
e bën këtë strukturë të dukshme. Për shkak se ka vetëm $m$ gjendje, vargu duhet të përsëritet; perioda dhe korrelacionet varen nga parametrat. Edhe kur histogrami i $u_n$ duket uniform, çiftet $(u_n,u_{n+1})$ mund të shtrihen në një numër të vogël hiperplanesh. Prandaj uniformiteti njëpërmasor nuk provon cilësinë shumëpërmasore.

Gjeneratorët modernë kanë gjendje shumë më të madhe dhe cilësi të testuar, por mbeten deterministë. Fara përcakton trajektoren. Ruajtja e farës siguron riprodhueshmëri, ndërsa përdorimi i së njëjtës farë për eksperimente që duhet të jenë të pavarura mund të krijojë korrelacione artificiale.

\subsection{Transformimi i anasjellë}
Nëse $U\sim\mathcal U(0,1)$ dhe $F$ është funksioni kumulativ i vazhdueshëm e rritës, atëherë
\begin{equation}
 X=F^{-1}(U)
\end{equation}
ka funksion kumulativ $F$, sepse
\begin{equation}
 P(X\le x)=P(F^{-1}(U)\le x)=P(U\le F(x))=F(x).
\end{equation}
Për eksponencialen $F(x)=1-e^{-\lambda x}$ merret
\begin{equation}
 X=-\frac{1}{\lambda}\ln(1-U).
\end{equation}
Shkrimi $\log1p(-U)$ është më i saktë kur $U$ është i vogël. Metoda është e thjeshtë kur inversi është i njohur; për shpërndarje të tjera nevojitet zgjidhje numerike ose metodë alternative.

\subsection{Box--Muller për mostrën Gauss}
Dy ndryshore Gauss standarde kanë dendësi të përbashkët
\begin{equation}
 f(z_1,z_2)=\frac{1}{2\pi}\exp[-(z_1^2+z_2^2)/2].
\end{equation}
Në koordinata polare, elementi i sipërfaqes është $r\dd r\dd\theta$, kështu që $\Theta$ është uniforme në $[0,2\pi)$ dhe $R^2$ ka shpërndarje eksponenciale me parametër $1/2$. Me $U_1,U_2$ uniforme të pavarura,
\begin{equation}
 R=\sqrt{-2\ln U_1},\qquad \Theta=2\pi U_2,
\end{equation}
dhe
\begin{equation}
 Z_1=R\cos\Theta,qquad Z_2=R\sin\Theta.
\end{equation}
Duhet shmangur $U_1=0$. Mostrat verifikohen jo vetëm me histogram, por edhe me mesatare, variancë, kuantile dhe korrelacion ndërmjet $Z_1$ e $Z_2$.

\subsection{Kampionimi me refuzim}
Le të jetë $f$ dendësia e synuar dhe $g$ dendësia propozuese, me $f(x)\le Mg(x)$. Gjenerohet $Y\sim g$ dhe $U\sim\mathcal U(0,1)$; $Y$ pranohet kur
\begin{equation}
 U\le\frac{f(Y)}{Mg(Y)}.
\end{equation}
Dendësia e një vlere të pranuar është proporcionale me
\begin{equation}
 g(y)\frac{f(y)}{Mg(y)}=\frac{f(y)}{M},
\end{equation}
pra pas normalizimit është $f$. Probabiliteti mesatar i pranimit është $1/M$. Një mbështjellëse e dobët e bën metodën të kushtueshme, sidomos në dimension të lartë.

\subsection{Kampionimi i rëndësisë dhe madhësia efektive}
Për një pritje nën $p$,
\begin{equation}
 I=\int h(x)p(x)\dd x
 =\int h(x)\frac{p(x)}{q(x)}q(x)\dd x,
\end{equation}
gjenerohen $X_i\sim q$ dhe përdoren peshat $w_i=p(X_i)/q(X_i)$. Kushti themelor është që $q>0$ kudo ku $|h|p$ kontribuon. Varianca mund të zvogëlohet shumë nëse $q$ vendos mostra në rajonet me kontribut të madh, por peshat shumë të pabarabarta e bëjnë vlerësimin të paqëndrueshëm. Një diagnostikë është
\begin{equation}
 N_{\mathrm{eff}}=\frac{(\sum_iw_i)^2}{\sum_iw_i^2}.
\end{equation}
Kjo nuk është identike me madhësinë efektive të një zinxhiri Markov, megjithëse të dyja matin humbjen e informacionit kundrejt $N$ mostrave ideale.

\subsection{Testimi i mostrave}
Një paketë minimale kontrollesh përfshin histogramin, CDF-në empirike, testin e momenteve dhe autokorrelacionin. Vlera-$p$ nuk provon rastësinë; ajo mat sa i pazakontë është statistiku nën një hipotezë të caktuar. Me shumë teste shfaqen rezultate të vogla edhe për një gjenerator të mirë, prandaj duhet parë modeli i përgjithshëm dhe jo vetëm minimumi i vlerave-$p$.

\subsection{Përmbledhje dhe ushtrime}
Gjenerimi i mostrës është pjesë e algoritmit shkencor. Një transformim i saktë matematikisht mund të dështojë numerikisht ose të trashëgojë korrelacionet e gjeneratorit bazë.
\begin{enumerate}
\item Gjeni kushtet për periodë të plotë të një gjeneratori kongruencial linear dhe testojini në parametra të vegjël.
\item Nxirrni transformimin e anasjellë për shpërndarjen Pareto.
\item Vërtetoni Box--Muller-in duke përfshirë Jakobinin polar.
\item Ndërtoni një mbështjellëse refuzimi për një dendësi trekëndore dhe parashikoni normën e pranimit.
\item Krahasoni dy propozime të kampionimit të rëndësisë me anë të variancës dhe $N_{\mathrm{eff}}$.
\end{enumerate}
